\documentclass[11pt,a4paper]{article}
\usepackage[utf8]{inputenc}
\usepackage[UTF8]{ctex} % 支持中文显示
\usepackage{amsmath, amssymb, amsfonts}
\usepackage{geometry}
\usepackage{bm}
\usepackage{tcolorbox}
% 在文档最开头添加以下宏包
\usepackage{booktabs}   % 解决表格 \toprule 等命令报错
\usepackage{amsmath}    % 解决数学公式环境报错
\usepackage{array}      % 增强表格功能
\usepackage[
    colorlinks=true,
    linkcolor=blue,
    unicode=true,          % 关键：使用 unicode 编码处理书签
    pdfpagemode=UseNone,   % 避免打开时自动弹出书签栏导致某些预览器崩溃
    pdfstartview=FitH
]{hyperref}
\geometry{left=2cm, right=2cm, top=2.5cm, bottom=2.5cm}

\title{YOLO 目标检测算法}
\author{魏舟}
\date{\today}


\begin{document}

\maketitle

\tableofcontents
\newpage

% --- 填充内容：第一部分 目标检测基础理论 ---

\section{目标检测基础理论}

\subsection{目标检测任务定义}
目标检测的核心目标是“在图像中找到感兴趣的目标，并确定其类别和位置”。
\begin{itemize}
    \item \textbf{分类 (Classification)}：识别图像中的物体属于哪个类别（如：猫、狗、汽车）。
    \item \textbf{定位 (Localization)}：确定物体在图像中的具体空间位置。
    \item \textbf{边界框 (Bounding Box)}：通常由四元组表示：
    \begin{itemize}
        \item 中心点坐标格式：$(x_{center}, y_{center}, w, h)$
        \item 左上角坐标格式：$(x_{min}, y_{min}, x_{max}, y_{max})$
    \end{itemize}
    \item \textbf{置信度 (Confidence Score)}：模型对其预测结果真实性的信心程度，计算公式通常为 $P(Object) \times \text{IoU}_{pred}^{truth}$。
\end{itemize}

\subsection{评价指标 (Metrics)}
为了客观衡量模型性能，通常使用以下指标：

\subsubsection{交并比 (IoU)}
交并比用于衡量预测框 $B_p$ 与真实框 $B_g$ 的重叠程度：
$$ \text{IoU} = \frac{\text{Area}(B_p \cap B_g)}{\text{Area}(B_p \cup B_g)} $$



\subsubsection{非极大值抑制 (NMS)}
在预测阶段，模型会对同一个物体生成多个重叠的候选框。**NMS** 的作用是：
\begin{enumerate}
    \item 按照置信度得分对所有候选框进行排序。
    \item 保留得分最高的框，并删除与之 IoU 大于预设阈值的所有其他框。
    \item 重复此过程直到所有框被处理。
\end{enumerate}

\subsubsection{mAP (mean Average Precision)}
这是衡量检测器综合性能的最重要指标：
\begin{itemize}
    \item \textbf{Precision (精确率)}：查准率，所有预测为正样本中实际为正的比例。
    \item \textbf{Recall (召回率)}：查全率，所有实际正样本中被预测出来的比例。
    \item \textbf{AP}：P-R 曲线下的面积。
    \item \textbf{mAP}：对所有类别的 AP 取算术平均。通常关注 mAP@0.5 或 mAP@.5:.95。
\end{itemize}

% --- 填充内容：第二部分 YOLO 核心思想 ---

\section{YOLO 核心思想：One-Stage Detector}

\subsection{设计哲学}
YOLO (You Only Look Once) 的核心理念是将目标检测重新定义为一个单一的\textbf{回归问题 (Regression Problem)}。
\begin{itemize}
    \item \textbf{统一网络}：从输入图像到边界框坐标和类别的概率，直接通过一个卷积神经网络完成推理。
    \item \textbf{网格划分 (Grid Cell)}：
    \begin{itemize}
        \item 将输入图像划分为 $S \times S$ 个网格。
        \item 如果物体的中心落在某个网格内，该网格就负责检测该物体。
        \item 每个网格预测 $B$ 个边界框及这些框的置信度，以及 $C$ 个类别的概率。
    \end{itemize}
    \item \textbf{全局信息}：由于使用全图信息进行训练和推理，YOLO 相比滑动窗口方法，背景误检（False Positives）更少。
\end{itemize}



\subsection{损失函数结构 (Loss Function)}
YOLOv1 使用均方误差 (MSE) 作为损失函数。为了平衡坐标定位、目标置信度和类别分类，损失函数由四部分加权组成：

\begin{equation}
\begin{aligned}
L = & \lambda_{coord} \sum_{i=0}^{S^2} \sum_{j=0}^{B} \mathbb{1}_{ij}^{obj} \left[ (x_i - \hat{x}_i)^2 + (y_i - \hat{y}_i)^2 + (\sqrt{w_i} - \sqrt{\hat{w}_i})^2 + (\sqrt{h_i} - \sqrt{\hat{h}_i})^2 \right] \\
& + \sum_{i=0}^{S^2} \sum_{j=0}^{B} \mathbb{1}_{ij}^{obj} (C_i - \hat{C}_i)^2 \\
& + \lambda_{noobj} \sum_{i=0}^{S^2} \sum_{j=0}^{B} \mathbb{1}_{ij}^{noobj} (C_i - \hat{C}_i)^2 \\
& + \sum_{i=0}^{S^2} \mathbb{1}_{i}^{obj} \sum_{c \in classes} (p_i(c) - \hat{p}_i(c))^2
\end{aligned}
\end{equation}

\subsubsection{关键点解析}
\begin{itemize}
    \item \textbf{权重系数}：$\lambda_{coord}$（通常取 5）用于增加坐标损失的权重；$\lambda_{noobj}$（通常取 0.5）用于惩罚不包含物体的背景框，防止负样本淹没正样本。
    \item \textbf{平方根处理}：对宽 $w$ 和高 $h$ 取平方根，是为了减小大尺度边界框对误差的敏感度（使小框的偏移比大框的相同像素偏移产生更大的 Loss）。
    \item \textbf{指示函数}：$\mathbb{1}_{ij}^{obj}$ 表示第 $i$ 个网格的第 $j$ 个预测框负责该物体。
\end{itemize}


% --- 填充内容：第三部分 YOLO 经典版本演进 ---

\section{YOLO 经典版本演进}

\subsection{YOLOv1：开创先河}
YOLOv1 证明了单阶段检测器在实时性上的巨大潜力。
\begin{itemize}
    \item \textbf{网络架构 (Backbone)}：参考了 GoogLeNet，包含 24 个卷积层和 2 个全连接层，输入分辨率为 $448 \times 448$。
    \item \textbf{局限性分析}：
    \begin{itemize}
        \item \textbf{空间约束}：每个网格只能预测一个类别，且只有两个边界框，导致其对密集目标（如：鸟群）检测效果较差。
        \item \textbf{泛化能力}：对物体的长宽比和尺度变化较为敏感。
        \item \textbf{定位误差}：主要的误差来源是定位不准（Localization Error）。
    \end{itemize}
\end{itemize}



\subsection{YOLOv2 (YOLO9000)：更好、更快、更强}
v2 版本引入了多项改进，使其在保持速度的同时大幅提升了精度。
\begin{itemize}
    \item \textbf{Batch Normalization}：每一层卷积后都加入 BN 层，收敛更快，并提升了 2\% 的 mAP。
    \item \textbf{Anchor Boxes 的先验知识}：
    \begin{itemize}
        \item 放弃了全连接层预测坐标，改用 Anchor Boxes。
        \item \textbf{Dimension Clusters}：使用 K-means 聚类算法在训练集中自动找到最合适的 Anchor 宽高比。
    \end{itemize}
    \item \textbf{Pass-through 层}：将高分辨率特征图（浅层）与低分辨率特征图（深层）拼接，增强了对小目标的检测能力。
    \item \textbf{多尺度训练}：每隔 10 个 batch 随机改变输入图像的尺寸（从 $320$ 到 $608$），增强模型的尺度适应性。
\end{itemize}

\subsection{YOLOv3：工业界的里程碑}
v3 是目前学术界和工业界参考最多的经典版本，结构非常稳健。
\begin{itemize}
    \item \textbf{Darknet-53 特征提取器}：引入了残差连接 (Residual Connections)，比 Darknet-19 更深，精度接近 ResNet-101 但速度更快。
    \item \textbf{多尺度预测 (FPN 思想)}：
    \begin{itemize}
        \item 在三个不同的尺度上进行预测（$13 \times 13$, $26 \times 26$, $52 \times 52$）。
        \item 借鉴了特征金字塔网络 (FPN) 的特征融合方式。
    \end{itemize}
    \item \textbf{逻辑回归替代 Softmax}：
    \begin{itemize}
        \item 使用多个独立逻辑分类器 (Binary Cross Entropy) 替代全类别的 Softmax。
        \item 支持了目标类别的层级关系（如：同一个物体既是“女人”又是“人”）。
    \end{itemize}
\end{itemize}

% --- 填充内容：第四部分 现代 YOLO 优化策略 ---

\section{现代 YOLO 优化策略 (v4 - v11+)}

\subsection{Backbone 增强}
主干网络的改进旨在以更低的计算成本提取更具表达力的特征。
\begin{itemize}
    \item \textbf{CSPNet (Cross Stage Partial Networks)}：
    \begin{itemize}
        \item 将特征图分为两部分，一部分经过密集块（Dense Block），另一部分直接与处理后的特征图合并。
        \item 解决了梯度信息重复的问题，在减少计算量的同时保持甚至提升了精度。
    \end{itemize}
    \item \textbf{RepVGG 与结构重参数化}：
    \begin{itemize}
        \item \textbf{训练时}：使用多分支结构（$3 \times 3$, $1 \times 1$ 及直连）来获取更好的特征表达。
        \item \textbf{推理时}：将所有分支融合为单一的 $3 \times 3$ 卷积层，实现“推理无损加速”。这是 YOLOv6 及以后版本的核心。
    \end{itemize}
\end{itemize}



\subsection{Neck 层结构}
Neck 层负责特征融合，连接 Backbone 与 Prediction Head。
\begin{itemize}
    \item \textbf{PANet (Path Aggregation Network)}：
    \begin{itemize}
        \item 在 FPN 的基础上增加了一个自底向上的路径（Bottom-up Path），使底层定位信息更容易传递到高层特征。
    \end{itemize}
    \item \textbf{SPP/SPPF (Spatial Pyramid Pooling - Fast)}：
    \begin{itemize}
        \item 通过不同维度的最大池化并拼接，极大地增加了感受野。
        \item \textbf{SPPF} 是 YOLOv5 提出的优化版本，将串行池化代替并行池化，速度大幅提升。
    \end{itemize}
\end{itemize}



\subsection{数据增强与训练技巧 (Bag of Freebies)}
这些技巧只增加训练成本，而不影响推理时间。
\begin{itemize}
    \item \textbf{Mosaic 数据增强}：
    \begin{itemize}
        \item 将四张图片拼接为一张，丰富了检测目标的上下文，变相增加了 Batch Size。
    \end{itemize}
    \item \textbf{Label Smoothing}：
    \begin{itemize}
        \item 防止模型对分类过于自信，提升泛化能力。
    \end{itemize}
    \item \textbf{IoU Loss 的演进}：
    \begin{itemize}
        \item \textbf{GIoU}：解决了 IoU 为 0 时梯度消失的问题。
        \item \textbf{DIoU}：考虑了边界框中心点的距离。
        \item \textbf{CIoU}：同时考虑了重叠面积、中心点距离和长宽比的对称性。
    \end{itemize}
\end{itemize}



\subsubsection{最新前沿 (v10/v11)}
\begin{itemize}
    \item \textbf{NMS-Free}：通过双重分配（Dual Assignment）消除对 NMS 的依赖（v10 核心）。
    \item \textbf{C3k2/C2f 结构}：针对不同尺度进行更精细的特征管理。
\end{itemize}


% --- 填充内容：第五部分 实战与部署总结 ---

\section{实战与部署总结}

\subsection{开发环境配置}
目前 YOLO 系列（尤其是 v5, v8, v10, v11）主要基于 Ultralytics 框架，其生态极大地简化了开发流程。
\begin{itemize}
    \item \textbf{框架选择}：
    \begin{itemize}
        \item \texttt{pip install ultralytics}：一键安装环境，集成训练、验证、预测。
        \item 硬件依赖：CUDA 工具包与 cuDNN 的版本匹配。
    \end{itemize}
    \item \textbf{数据集准备}：
    \begin{itemize}
        \item \textbf{YOLO 格式}：每张图对应一个 \texttt{.txt} 文件，内容为 \texttt{<class-id> <x> <y> <w> <h>}（归一化坐标）。
        \item \textbf{格式转换}：
            \begin{itemize}
                \item \textbf{Pascal VOC}：XML 标签格式。
                \item \textbf{MS COCO}：JSON 嵌套格式。
            \end{itemize}
        \item 使用 \texttt{data.yaml} 定义数据集路径和类别名称。
    \end{itemize}
\end{itemize}

\subsection{模型量化与加速}
为了满足嵌入式设备或高并发服务器的实时性要求，通常需要对模型进行转换和压缩。

\subsubsection{ONNX 导出注意事项}
ONNX 是模型部署的中间媒介。导出时需注意：
\begin{itemize}
    \item \textbf{动态维度 (Dynamic Axes)}：是否允许输入图片尺寸可变。
    \item \textbf{算子兼容性}：部分 YOLO 的自定义层（如某些版本的 Upsample）在导出时可能需要调整 \texttt{opset\_version}。
    \item \textbf{后处理集成}：是否将 NMS 逻辑包含在 ONNX 模型内部。
\end{itemize}



\subsubsection{TensorRT 部署流程}
TensorRT 是 NVIDIA 推出的高性能推理引擎，其核心在于对模型进行重构。
\begin{itemize}
    \item \textbf{层融合 (Layer Fusion)}：将卷积、偏移和激活函数融合为一个核函数（Kernel）。
    \item \textbf{量化 (Quantization)}：
        \begin{itemize}
            \item \textbf{FP16}：在精度几乎无损的情况下，推理速度显著提升。
            \item \textbf{INT8 (PTQ/QAT)}：通过校准集（Calibration set）将权重压缩至 8 位，极大地减少内存占用，但需注意精度回退风险。
    \end{itemize}
    \item \textbf{推理引擎生成}：在目标设备上生成专用的 \texttt{.engine} 文件，以发挥最大算力。
\end{itemize}

\begin{table}[h]
    \centering
    \caption{常见部署格式对比}
    \begin{tabular}{@{}llll@{}}
        \toprule
        \textbf{格式} & \textbf{平台} & \textbf{优势} & \textbf{适用场景} \\ \midrule
        PyTorch (.pt) & GPU & 灵活、易于调试 & 实验、训练阶段 \\
        ONNX (.onnx) & 跨平台 & 通用性强 & 中间转换、CPU 推理 \\
        TensorRT (.engine) & NVIDIA GPU & 极致性能 & 边缘计算、实时视频流 \\
        TFLite (.tflite) & 移动端/端侧 & 轻量化 & Android/iOS 应用 \\ \bottomrule
    \end{tabular}
\end{table}

% --- 填充内容：结语与参考文献 ---

\section{结语与展望}
随着深度学习技术的飞速发展，YOLO 已经从最初的纯卷积架构演变为一个高度集成的视觉感知框架。
\begin{itemize}
    \item \textbf{端到端检测 (NMS-free)}：YOLOv10 引入的双重分配机制（Dual Assignment）使得模型不再依赖后处理 NMS，显著降低了推理延迟。
    \item \textbf{多模态与开放词汇 (OVD)}：未来的 YOLO 将不仅限于检测预定义的 80 个类别，通过结合 CLIP 等大模型，实现“用自然语言描述什么就检测什么”。
    \item \textbf{实时实例分割}：如 YOLOv8/v11-seg，通过增加计算轻量级的 Proto 模块，使实时像素级理解成为可能。
\end{itemize}



\section{参考文献与资源}

\subsection{核心论文列表}
\begin{itemize}
    \item \textbf{YOLOv1}: Redmon, et al. "You Only Look Once: Unified, Real-Time Object Detection." CVPR 2016.
    \item \textbf{YOLOv3}: Redmon, et al. "YOLOv3: An Incremental Improvement." arXiv 2018.
    \item \textbf{YOLOv4}: Bochkovskiy, et al. "YOLOv4: Optimal Speed and Accuracy of Object Detection." arXiv 2020.
    \item \textbf{YOLOv7}: Wang, et al. "YOLOv7: Trainable Bag-of-freebies Sets New State-of-the-art for Real-time Object Detectors." CVPR 2023.
    \item \textbf{YOLOv10}: Wang, et al. "YOLOv10: Real-Time End-to-End Object Detection." arXiv 2024.
\end{itemize}

\subsection{优秀开源库与工具}
\begin{itemize}
    \item \textbf{Ultralytics (YOLOv8/v11)}: \url{https://github.com/ultralytics/ultralytics} —— 目前最易用的生产力工具。
    \item \textbf{MMDetection}: \url{https://github.com/open-mmlab/mmdetection} —— 覆盖几乎所有主流检测算法的学术级框架。
    \item \textbf{Roboflow}: 用于在线标注数据集并导出 YOLO 格式的利器。
\end{itemize}

\subsection{学习建议}
\begin{quote}
    “不仅要看模型跑出的 mAP 分数，更要深入阅读其 \texttt{head.py} 和 \texttt{loss.py} 的源码。理解样本匹配（Assigner）的逻辑，比理解网络架构本身更重要。”
\end{quote}
\end{document}